\documentclass[10pt,a4paper]{amsart}
\usepackage[margin=19mm]{geometry}
\usepackage{amsmath,amssymb,mathtools}
\usepackage[scr=rsfs,cal=euler]{mathalfa}
\usepackage{xfrac}
\usepackage[unicode,hidelinks,bookmarks=false]{hyperref}
\newcommand{\ie}{i.e.\ }
\newcommand{\cf}{cf.\ }
\newcommand{\resp}{respectively\ }
\newcommand{\Subsection}{\subsection}
\providecommand{\nobreakdash}{\nobreak\hbox{-}}
\setlength{\emergencystretch}{2em}
\multlinegap0pt
\usepackage{amssymb}
\usepackage{enumerate}
\usepackage{float}
\usepackage[shortlabels]{enumitem}
\usepackage{xcolor}
\newtheorem{proposition}{Proposition}
\title{Hom-unitality and hom-associative structures}
\author{Germ\'an Garc\'ia Butenegro, Abdennour Kitouni, Sergei Silvestrov}
\date{Source: arXiv:2602.07314v2 (CC BY 4.0)}
\begin{document}
\maketitle

\noindent\emph{Source and licence.} Hom-unitality and hom-associative structures. Version arXiv:2602.07314v2 (CC BY 4.0), distributed by arXiv under the Creative Commons Attribution 4.0 International licence, http://creativecommons.org/licenses/by/4.0/, as recorded in the arXiv licence metadata for this exact version and shown on the abstract page https://arxiv.org/abs/2602.07314v2. Full licence notice: \texttt{licenses/arxiv-2602.07314-CC-BY-4.0.txt}.\par\smallskip

\noindent\textit{Editorial scope.} Counted unit: the article's proposition prop:LeibHomUnities (source0.tex lines 1737--1743). The article's complete original proof of it is delivered with it (source0.tex lines 1745--1783, one \texttt{proof} environment of 39 lines). No mathematics is written, repaired or re-derived: the statement and the proof are the article's own lines, in the article's own notation, and no retained line is substituted or re-flowed.  No further source-local context is needed: every cross-reference of every retained line resolves inside the two delivered spans.  Nothing was omitted from any retained span, so the package carries no omission record.  No retained line contains a citation, so the wrapper carries no bibliography and no bibliography entry.  No retained line contains a figure environment, an image-inclusion command or drawing code, so no image is delivered and the package declares no asset dependency.  Editorial formatting only, in addition: an AMS \texttt{amsart} preamble that restates the article's own declarations and macros used by the retained lines, namely the environments \texttt{proposition} on separate counters, together with the standard packages those lines need.  The retained environments print in this document as Proposition 1 in order of appearance, whereas the article prints the same items with the numbering of its own full document.  The complete native source of the article as distributed in the arXiv e-print for this exact version is bundled unchanged as \texttt{sources/194171/source0.tex}.
            \begin{proposition}\label{prop:LeibHomUnities}
				Let $(\mathcal{L},[\cdot,\cdot])$ be a Leibniz algebra. Then $(\mathcal{L},[\cdot,\cdot],L_a)$ is hom-associative for all $a$ in the subspace
				\begin{equation*}
					HU_n(\mathcal{L})=C(\mathcal{L})\cap Ann_\mathcal{L}^l([\mathcal{L},\mathcal{L}])=C(\mathcal{L})\cap Ann_\mathcal{L}^r([\mathcal{L},\mathcal{L}])=C(\mathcal{L})\cap Ann_\mathcal{L}([\mathcal{L},\mathcal{L}]).
				\end{equation*}
				Moreover, every element in $HU_n(\mathcal{L})$ is 3-nilpotent.
			\end{proposition}

			\begin{proof}
				Compute the associator of three elements $a,b,c\in\mathcal{L}$ using the Leibniz rule:
				\begin{align*}
					\textbf{Left: }[a,[b,c]]=[[a,b],c]+[b,[a,c]]\Rightarrow [a,b,c]_{\mathrm{as}}=-[b,[a,c]].\\
					\textbf{Right: }[[a,b],c]=[[a,c],b]+[a,[b,c]]\Rightarrow [a,b,c]_{\mathrm{as}}=[[a,c],b].
				\end{align*}
				It follows that $[\mathcal{L},\mathcal{L},\mathcal{L}]_{\mathrm{as}}=[[\mathcal{L},\mathcal{L}],\mathcal{L}]$ (left) and $[\mathcal{L},\mathcal{L},\mathcal{L}]_{\mathrm{as}}=[\mathcal{L},[\mathcal{L},\mathcal{L}]]$ (right).
				
				Any $a\in N(\mathcal{L})$ satisfies, for all $b,c\in\mathcal{L}$, the following relations if $\mathcal{L}$ is left Leibniz:
				\begin{align*}
					\left.\begin{array}{c}
						0=[c,b,a]_{\mathrm{as}}=[[c,b],a]-[c,[b,a]]=-[b,[c,a]]  \\
						0=[b,c,a]_{\mathrm{as}}=[[b,c],a]-[b,[c,a]]=-[c,[b,a]]
					\end{array}\right\}&\Rightarrow\ \ [[b,c],a]=[c,[b,a]]=0\\
					\left.\begin{array}{c}
						0=[a,b,c]_{\mathrm{as}}=[[a,b],c]-[a,[b,c]]=-[b,[a,c]]  \\
						0=[b,a,c]_{\mathrm{as}}=[[b,a],c]-[b,[a,c]]=-[c,[a,b]]  \\
					\end{array}\right\}&\Rightarrow \left\{\begin{matrix}
						[[a,b],c]=[[b,a],c]=0\\
						[a,[b,c]]=[b,[a,c]]=0
					\end{matrix} \right.
				\end{align*}
				and the following relations if $\mathcal{L}$ is right-Leibniz:
				\begin{align*}
					\left.\begin{array}{c}
						0=[a,b,c]_{\mathrm{as}}=[[a,b],c]-[a,[b,c]]=[[a,c],b]  \\
						0=[a,c,b]_{\mathrm{as}}=[[a,c],b]-[a,[c,b]]=[[a,b],c]
					\end{array}\right\}&\Rightarrow\ \ [a,[b,c]]=[[a,c],b]=0\\
					\left.\begin{array}{c}
						0=[c,b,a]_{\mathrm{as}}=[[c,b],a]-[c,[b,a]]=[[c,a],b]  \\
						0=[c,a,b]_{\mathrm{as}}=[[c,a],b]-[c,[a,b]]=[[c,b],a]
					\end{array}\right\}&\Rightarrow \left\{\begin{matrix}
						[c,[b,a]]=[c,[a,b]]=0\\
						[[c,b],a]=[[c,a],b]=0
					\end{matrix} \right.
				\end{align*}
				Every triple product involving $a$ is zero. In particular, $a$ commutes with any $[b,c]$ and is 3-nilpotent. That is, $N(\mathcal{L})$, and in particular $C(\mathcal{L})\cap N(\mathcal{L})$, is a subspace of the subspace of 3-nilpotent elements of $\mathcal{L}$. By the Lebniz identity,
				$[\mathcal{L},\mathcal{L},\mathcal{L}]_{\mathrm{as}}=[[\mathcal{L},\mathcal{L}],\mathcal{L}]\subseteq [\mathcal{L},\mathcal{L}]$ and thus  $a\in N(\mathcal{L})\Rightarrow a\in Ann_\mathcal{L}^l([\mathcal{L},\mathcal{L}])\subseteq Ann_\mathcal{L}^l([[\mathcal{L},\mathcal{L}],\mathcal{L}])= Ann_\mathcal{L}^l([\mathcal{L},\mathcal{L},\mathcal{L}]_{\mathrm{as}})$. This indicates that $C(\mathcal{L})\cap N(\mathcal{L})$ cancels $[\mathcal{L},\mathcal{L}]$ on both sides by centrality, so we can use left, right, or two-sided annihilator notation indistinctly. This completes the proof.
			\end{proof}				

\end{document}
